\section{Repaso de funciones}

Una función $f\colon A \to B$ consta de tres partes: un conjunto $A$ (dominio), un conjunto $B$ (codominio) y una regla que asocia a cada elemento $x \in A$ un único elemento $f(x) \in B$.

El \emph{gráfico} de $f$ es el subconjunto de $A \times B$ dado por
\[
  G(f) = \{ (x, y) \in A \times B : y = f(x) \}.
\]

\begin{definition}[Inyectividad]
  $f\colon A \to B$ es \emph{inyectiva} si para todo $x, y \in A$ vale $f(x) = f(y) \Rightarrow x = y$.
\end{definition}

\begin{definition}[Suryectividad]
  $f\colon A \to B$ es \emph{suryectiva} si para todo $y \in B$ existe $x \in A$ tal que $f(x) = y$.
\end{definition}

\begin{definition}[Biyectividad]
  $f\colon A \to B$ es \emph{biyectiva} si es inyectiva y suryectiva.
\end{definition}

\begin{definition}[Imagen]
  Dados $f\colon A \to B$ y $X \subset A$, la \emph{imagen} de $X$ por $f$ es $f(X) = \{ f(x) : x \in X \}$.
\end{definition}

\subsection{Propiedades de la imagen}

\begin{prop}
  Sean $X, Y \subset A$. Entonces:
  \begin{enumerate}[label=(\arabic*)]
    \item $f(X \cup Y) = f(X) \cup f(Y)$.
    \item $X \subset Y \Rightarrow f(X) \subset f(Y)$.
    \item $f(\varnothing) = \varnothing$.
  \end{enumerate}
\end{prop}

\begin{prop}
  Si $X, Y \subset A$, entonces $f(X \cap Y) \subset f(X) \cap f(Y)$. La igualdad vale para todo par $X, Y$ si y solo si $f$ es inyectiva.
\end{prop}

\begin{proof}
  Si $b \in f(X \cap Y)$, existe $x \in X \cap Y$ con $f(x) = b$. Como $x \in X$ e $x \in Y$, resulta $b \in f(X) \cap f(Y)$.

  Si $f$ es inyectiva y $b \in f(X) \cap f(Y)$, existen $a \in X$ y $a' \in Y$ con $f(a) = f(a') = b$. Por inyectividad $a = a'$, luego $a \in X \cap Y$ y $b \in f(X \cap Y)$.

  Si $f$ no es inyectiva, existen $x \neq y$ con $f(x) = f(y)$. Tomando $X = \{x\}$ e $Y = \{y\}$ se tiene $X \cap Y = \varnothing$, de modo que $f(X \cap Y) = \varnothing$, pero $f(X) \cap f(Y) = \{f(x)\} \neq \varnothing$.
\end{proof}

\subsection{Preimagen}

\begin{definition}[Preimagen]
  Dados $f\colon A \to B$ e $Y \subset B$, la \emph{preimagen} de $Y$ por $f$ es
  \[
    f^{-1}(Y) = \{ x \in A : f(x) \in Y \}.
  \]
\end{definition}

\begin{prop}
  Sean $X, Y \subset B$. Entonces:
  \begin{enumerate}[label=(\arabic*)]
    \item $f^{-1}(X \cup Y) = f^{-1}(X) \cup f^{-1}(Y)$.
    \item $f^{-1}(X \cap Y) = f^{-1}(X) \cap f^{-1}(Y)$.
    \item $f^{-1}(Y^c) = {\bigl(f^{-1}(Y)\bigr)}^c$.
    \item $X \subset Y \Rightarrow f^{-1}(X) \subset f^{-1}(Y)$.
    \item $f^{-1}(B) = A$ y $f^{-1}(\varnothing) = \varnothing$.
  \end{enumerate}
\end{prop}

\begin{proof}
  Todas se deducen directamente de la definición. Por ejemplo, para (3):
  \[
    x \in f^{-1}(Y^c) \iff f(x) \notin Y \iff x \notin f^{-1}(Y) \iff x \in {\bigl(f^{-1}(Y)\bigr)}^c. \qedhere
  \]
\end{proof}

\subsection{Composición e inversas}

Sean $f\colon A \to B$ y $g\colon B \to C$. La \emph{composición} $g \circ f\colon A \to C$ se define como $(g \circ f)(x) = g(f(x))$ para todo $x \in A$. Basta con que $f(A)$ esté contenido en el dominio de $g$.

\begin{prop}
  La composición de funciones inyectivas es inyectiva y la composición de funciones suryectivas es suryectiva. Además, ${(g \circ f)}^{-1}(Z) = f^{-1}\bigl(g^{-1}(Z)\bigr)$ para todo $Z \subset C$.
\end{prop}

\begin{definition}[Restricción]
  La \emph{restricción} de $f\colon A \to B$ a un subconjunto $X \subset A$ es la función $f|_X\colon X \to B$ dada por $f|_X(x) = f(x)$.
\end{definition}

\begin{definition}[Inversas laterales]
  Sea $f\colon A \to B$. Una función $g\colon B \to A$ es:
  \begin{itemize}
    \item una \emph{inversa a izquierda} de $f$ si $g \circ f = \id_A$;
    \item una \emph{inversa a derecha} de $f$ si $f \circ g = \id_B$;
    \item una \emph{inversa} de $f$ si es inversa a izquierda y a derecha.
  \end{itemize}
\end{definition}

\begin{prop}
  Sea $f\colon A \to B$ con $A \neq \varnothing$. Entonces:
  \begin{enumerate}[label=(\arabic*)]
    \item $f$ tiene inversa a izquierda si y solo si $f$ es inyectiva;
    \item $f$ tiene inversa a derecha si y solo si $f$ es suryectiva (esta implicación usa el axioma de elección);
    \item $f$ es biyectiva si y solo si tiene una inversa, que en tal caso es única.
  \end{enumerate}
\end{prop}

\subsection{Familias}

Sea $L$ un conjunto de \emph{índices}, cuyos elementos se representan genéricamente con $\lambda$. Dado un conjunto $X$, una \emph{familia} de elementos de $X$ con índices en $L$ es una función $x\colon L \to X$. El valor de $x$ en $\lambda \in L$ se nota $x_\lambda$ y la familia, ${(x_\lambda)}_{\lambda \in L}$.

\begin{eg}
  Si $L = \{1, 2\}$, una familia $x\colon \{1, 2\} \to X$ se representa por los valores $x_1$ y $x_2$, es decir, se identifica con un par ordenado $(x_1, x_2)$ de elementos de $X$.
\end{eg}

Una familia indexada por $\N$ se llama \emph{sucesión}: ${(x_n)}_{n \in \N}$ es una función $x\colon \N \to X$ con $x_n = x(n)$.

Las propiedades anteriores se extienden a familias arbitrarias de conjuntos; por ejemplo, $f\bigl(\bigcup_\lambda X_\lambda\bigr) = \bigcup_\lambda f(X_\lambda)$ y $f^{-1}\bigl(\bigcap_\lambda Y_\lambda\bigr) = \bigcap_\lambda f^{-1}(Y_\lambda)$.

\section{Números naturales}

Partimos de un conjunto $\N$ y una función $S\colon \N \to \N$ (el \emph{sucesor}) que cumple los \emph{axiomas de Peano}:
\begin{enumerate}[label=(\arabic*)]
  \item $S$ es inyectiva.
  \item $\N \setminus S(\N)$ tiene un único elemento, que llamamos $1$.
  \item (\emph{Principio de inducción}) Si $X \subset \N$ cumple $1 \in X$ y $S(m) \in X$ para todo $m \in X$, entonces $X = \N$.
\end{enumerate}

El principio de inducción permite definir operaciones recursivamente. La \emph{suma} se define por
\[
  m + 1 = S(m), \qquad m + S(n) = S(m + n).
\]

\begin{prop}[Asociatividad]
  Para todos $m, n, p \in \N$, $m + (n + p) = (m + n) + p$.
\end{prop}

\begin{proof}
  Sea $X = \{ p \in \N : m + (n + p) = (m + n) + p \text{ para todos } m, n \in \N \}$.

  Para $p = 1$: $m + (n + 1) = m + S(n) = S(m + n) = (m + n) + 1$, así que $1 \in X$.

  Si $p \in X$, entonces
  \[
    m + (n + S(p)) = m + S(n + p) = S\bigl(m + (n + p)\bigr) = S\bigl((m + n) + p\bigr) = (m + n) + S(p),
  \]
  de modo que $S(p) \in X$. Por inducción, $X = \N$.
\end{proof}

\begin{prop}
  La suma es conmutativa: $m + n = n + m$. Además vale la ley de cancelación: $m + n = m + p \Rightarrow n = p$.
\end{prop}

\begin{prop}[Tricotomía]
  Dados $m, n \in \N$, se cumple exactamente una de las siguientes condiciones:
  \[
    m = n, \qquad \exists\, p \in \N \colon m = n + p, \qquad \exists\, p \in \N \colon n = m + p.
  \]
  Se escribe $n < m$ en el segundo caso y $m < n$ en el tercero.
\end{prop}

\begin{definition}[Producto]
  La \emph{multiplicación} se define recursivamente por $m \cdot 1 = m$ y $m \cdot S(n) = m \cdot n + m$. Es asociativa, conmutativa, distributiva respecto de la suma, cumple la ley de cancelación y es compatible con el orden.
\end{definition}

\begin{theorem}[Principio de buena ordenación]
  Todo subconjunto no vacío $A \subset \N$ tiene un elemento mínimo.
\end{theorem}

\begin{proof}
  Para cada $m \in \N$ sea $I_m = \{ p \in \N : p \leq m \}$ y sea
  \[
    X = \{ m \in \N : I_m \subset \N \setminus A \}.
  \]
  Si $1 \in A$, entonces $1$ es el mínimo de $A$. Supongamos $1 \notin A$; entonces $1 \in X$. Además $X \neq \N$, porque si fuera $X = \N$ tendríamos $m \in I_m \subset \N \setminus A$ para todo $m$, es decir $A = \varnothing$.

  Por el principio de inducción, existe $m \in X$ tal que $m + 1 \notin X$. Entonces todos los elementos menores o iguales que $m$ están en $\N \setminus A$, pero $I_{m+1} \not\subset \N \setminus A$, de donde $m + 1 \in A$. Luego $m + 1$ es el mínimo de $A$.
\end{proof}

\section{Cuerpos}

Un \emph{cuerpo} es un conjunto $\K$ dotado de dos operaciones, suma y producto, que cumplen los siguientes axiomas para todos $x, y, z \in \K$.

\subsection{Axiomas de la suma}

\begin{enumerate}[label=(S\arabic*)]
  \item Asociatividad: $(x + y) + z = x + (y + z)$.
  \item Conmutatividad: $x + y = y + x$.
  \item Neutro: existe $0 \in \K$ tal que $x + 0 = x$.
  \item Opuesto: existe $-x \in \K$ tal que $x + (-x) = 0$.
\end{enumerate}

\subsection{Axiomas del producto}

\begin{enumerate}[label=(P\arabic*)]
  \item Asociatividad: $(xy)z = x(yz)$.
  \item Conmutatividad: $xy = yx$.
  \item Neutro: existe $1 \in \K$, $1 \neq 0$, tal que $x \cdot 1 = x$.
  \item Inverso: si $x \neq 0$, existe $x^{-1} \in \K$ tal que $x \cdot x^{-1} = 1$.
\end{enumerate}

Además se pide la \emph{distributividad}: $x(y + z) = xy + xz$.

\begin{eg}
  $\Q$, $\R$, $\C$ y $\Z_2 = \Z/2\Z$ son cuerpos.
\end{eg}

\subsection{Cuerpos ordenados}

Un \emph{cuerpo ordenado} es un cuerpo $\K$ con un subconjunto $P \subset \K$, llamado conjunto de elementos \emph{positivos}, tal que:
\begin{enumerate}[label=(O\arabic*)]
  \item si $x, y \in P$, entonces $x + y \in P$ y $xy \in P$;
  \item para todo $x \in \K$ se cumple exactamente una de las siguientes: $x \in P$, $-x \in P$ o $x = 0$.
\end{enumerate}

\begin{eg}
  $\Q$ con $P = \{ p/q : p, q \in \N \}$ es un cuerpo ordenado. En cambio, $\Z_2$ no admite un orden: si $1 \in P$, entonces $1 + 1 = 0 \in P$, lo que contradice (O2).
\end{eg}

\begin{prop}
  En un cuerpo ordenado, si $a \neq 0$ entonces $a^2 \in P$.
\end{prop}

\begin{proof}
  Si $a \in P$, entonces $a^2 = a \cdot a \in P$ por (O1). Si $a \notin P$, por (O2) es $-a \in P$ y entonces $a^2 = (-a)(-a) \in P$.
\end{proof}

En un cuerpo ordenado se define $x < y$ cuando $y - x \in P$. Esta relación cumple:

\begin{prop}
  Sean $x, y, z \in \K$. Entonces:
  \begin{enumerate}[label=(\arabic*)]
    \item (Transitividad) $x < y$ e $y < z$ implican $x < z$.
    \item (Tricotomía) vale exactamente una de las relaciones $x = y$, $x < y$, $y < x$.
    \item (Monotonía de la suma) $x < y \Rightarrow x + z < y + z$.
    \item (Monotonía del producto) $x < y$ y $0 < z \Rightarrow xz < yz$.
  \end{enumerate}
\end{prop}

Se escribe $x \leq y$ si $x < y$ o $x = y$, es decir, si $y - x \in P \cup \{0\}$. Con esta relación se cumplen todas las propiedades anteriores y además la antisimetría: $x \leq y$ e $y \leq x$ implican $x = y$.

También se tiene la noción de intervalo. Dados $a, b \in \K$ se definen
\[
  [a, b] = \{ x \in \K : a \leq x \leq b \}, \quad (a, b) = \{ x \in \K : a < x < b \},
\]
y análogamente $[a, b)$, $(a, b]$, $[a, +\infty)$, $(a, +\infty)$, $(-\infty, b]$ y $(-\infty, b)$. Notar que $[a, a] = \{a\}$.

Un subconjunto $X \subset \K$ es \emph{acotado superiormente} si existe $b \in \K$ tal que $x \leq b$ para todo $x \in X$ (se dice que $b$ es una \emph{cota superior}). Análogamente se define \emph{acotado inferiormente} y \emph{cota inferior}. Un conjunto es \emph{acotado} si lo es superior e inferiormente.
